\documentclass[11pt]{article}

\usepackage[margin=1in]{geometry}
\usepackage{amsmath, amsthm, amssymb, amsfonts}
\usepackage{mathtools}
\usepackage{booktabs}
\usepackage{enumitem}
\usepackage{microtype}
\usepackage{natbib}
\usepackage{xcolor}
\usepackage{hyperref}
\hypersetup{colorlinks=true, linkcolor=blue!50!black, citecolor=blue!50!black, urlcolor=blue!50!black}

\theoremstyle{plain}
\newtheorem{theorem}{Theorem}[section]
\newtheorem{lemma}[theorem]{Lemma}
\newtheorem{proposition}[theorem]{Proposition}
\newtheorem{corollary}[theorem]{Corollary}
\theoremstyle{definition}
\newtheorem{definition}[theorem]{Definition}
\newtheorem{assumption}[theorem]{Assumption}
\theoremstyle{remark}
\newtheorem{remark}[theorem]{Remark}

\newcommand{\R}{\mathbb{R}}
\newcommand{\C}{\mathbb{C}}
\newcommand{\E}{\mathbb{E}}
\newcommand{\norm}[1]{\left\lVert#1\right\rVert}
\newcommand{\abs}[1]{\left\lvert#1\right\rvert}
\newcommand{\spec}{\operatorname{spec}}
\newcommand{\trace}{\operatorname{tr}}
\DeclareMathOperator*{\argmin}{arg\,min}
\DeclareMathOperator*{\argmax}{arg\,max}

\title{\textbf{Draft Preliminaries}\\
\large Polynomial Input Preconditioning for Universal Time Series Forecasting}
\author{}
\date{}

\begin{document}
\maketitle

\paragraph{Purpose of this draft.}
This file is a proposed replacement for the mathematical preliminaries in the Junior Paper. It follows the organization of \emph{Intro to Online Control}~\citep{hazan2025onlinecontrol}: first define passive prediction in linear dynamical systems, then explain stability and finite-memory predictors, then state the online-learning guarantees for those predictors, then introduce spectral filtering as the bridge to universal sequence preconditioning. The final subsection connects that lineage to patch-based transformer forecasting.

%==============================================================================
\section{Preliminaries}
\label{sec:preliminaries-draft}
%==============================================================================

The mathematical background for our method has two roles. First, it explains why long-memory sequence prediction is easy for strictly stable systems but hard near the unit circle. Second, it identifies the specific object we borrow from the theory: a fixed polynomial filter that reshapes the spectrum of a sequence before it is passed to a predictor. We present only the statements needed for the paper; exact constants and full proofs are deferred to the cited sources.

%------------------------------------------------------------------------------
\subsection{Passive Prediction in Linear Dynamical Systems}
\label{sec:passive-lds}
%------------------------------------------------------------------------------

We use the linear-dynamical-system notation of \citet[Chs.~4 and~9]{hazan2025onlinecontrol}. A partially observed linear time-invariant system has hidden state $x_t \in \R^{d_x}$, observed input $u_t \in \R^{d_u}$, and output $y_t \in \R^{d_y}$:
\begin{equation}
    x_{t+1} = A x_t + B u_t + w_t,
    \qquad
    y_t = C x_t + v_t .
    \label{eq:lds-input-output}
\end{equation}
Here $A,B,C$ are unknown system matrices, while $w_t$ and $v_t$ denote process and observation noise. In time-series forecasting there may be no exogenous input; this is recovered by setting $u_t \equiv 0$ or by treating past observations as the relevant input history.

The prediction problem is passive: unlike control, the learner does not choose $u_t$ to influence the state. At each time $t$, the learner observes the past history and predicts $\hat y_t$; after observing $y_t$, it suffers loss $\ell_t(\hat y_t,y_t)$, usually squared error. For a predictor class $\Pi$, regret is
\begin{equation}
    \operatorname{Regret}_T(\mathcal A,\Pi)
    :=
    \sum_{t=1}^{T} \ell_t(\hat y_t,y_t)
    -
    \min_{\pi \in \Pi}
    \sum_{t=1}^{T} \ell_t(\hat y_t^\pi,y_t).
    \label{eq:prediction-regret}
\end{equation}
Sublinear regret means that the learner's average loss approaches that of the best predictor in the reference class. This is the sense in which the online-learning results below should be read.

%------------------------------------------------------------------------------
\subsection{Stability and Finite Memory}
\label{sec:stability-memory}
%------------------------------------------------------------------------------

The spectral radius of a square matrix is
\[
    \rho(A) := \max\{\abs{\lambda}: \lambda \in \spec(A)\}.
\]
The basic stability fact is that, for a time-invariant linear system, $\rho(A)<1$ is equivalent to $A^t \to 0$ as $t \to \infty$~\citep[Fact~4.2]{hazan2025onlinecontrol}. For finite-time learning bounds one usually needs the quantitative version
\begin{equation}
    \norm{A^t} \le M \rho^t
    \qquad \text{for all } t \ge 0,
    \qquad \rho \in (0,1),
    \label{eq:quantitative-stability}
\end{equation}
or, equivalently, $\norm{A^t} \le \kappa(1-\delta)^t$ with stability margin $\delta>0$. This says that the influence of the remote past decays geometrically.

The following truncation lemma is the precise version of the common statement that strictly stable systems have finite effective memory.

\begin{lemma}[Stable LDSs are approximated by finite-memory predictors; cf.\ {\citet[Lemma~9.6]{hazan2025onlinecontrol}}]
\label{lem:stable-truncation}
Consider the noiseless input-output LDS $x_{t+1}=Ax_t+Bu_t$, $y_t=Cx_t$. Suppose
\[
    \norm{A^t} \le M\rho^t,\qquad
    \norm{B}\le \beta,\qquad
    \norm{C}\le \gamma,\qquad
    \norm{u_t}\le U,
\]
with $\rho \in (0,1)$. Let $\pi_{A,B,C,x_0}$ denote the exact open-loop LDS predictor and let $\pi_{A,B,C,h}$ be the length-$h$ truncated predictor
\[
    \hat y_t^{(h)}
    =
    \sum_{i=1}^{h} C A^{i-1}B u_{t-i}.
\]
Then
\begin{equation}
    \sum_{t=1}^{T}
    \norm{\hat y_t^{\pi_{A,B,C,x_0}}-\hat y_t^{(h)}}
    \le
    T \cdot \frac{M\beta\gamma U}{1-\rho}\rho^h
    +
    \sum_{t=1}^{\min\{T,d_x\}} \norm{CA^t x_0}.
    \label{eq:stable-truncation}
\end{equation}
\end{lemma}

Thus, after the initial transient, a stable system can be approximated to error $\varepsilon$ using memory
\[
    h = O\!\left(\frac{\log(T/\varepsilon)}{1-\rho}\right)
\]
when $\rho$ is close to one. This is the useful ``strictly stable'' baseline: finite-memory linear prediction is valid, but the required memory can grow poorly as the stability margin $1-\rho$ shrinks.

%------------------------------------------------------------------------------
\subsection{Linear Predictors and Online Convex Optimization}
\label{sec:linear-predictors}
%------------------------------------------------------------------------------

The standard convex comparator class for passive prediction is the class of finite-memory linear predictors:
\begin{definition}[Finite-memory linear predictors]
\label{def:linear-predictors}
For memory lengths $h,k$ and norm radius $\kappa$, define $\Pi^L_{h,k,\kappa}$ to be the set of predictors
\begin{equation}
    \hat y_t
    =
    \sum_{i=1}^{h} M_{1,i} y_{t-i}
    +
    \sum_{j=1}^{k} M_{2,j} u_{t-j},
    \qquad
    \sum_{i=1}^{h}\norm{M_{1,i}}
    +
    \sum_{j=1}^{k}\norm{M_{2,j}}
    \le \kappa.
    \label{eq:linear-predictor}
\end{equation}
\end{definition}

This class is computationally attractive because it is convex in the predictor matrices. Online Gradient Descent (OGD) therefore gives an immediate learning guarantee.

\begin{theorem}[OGD for finite-memory predictors; cf.\ {\citet[Thm.~6.1 and Cor.~9.7]{hazan2025onlinecontrol}}]
\label{thm:ogd-linear-predictors}
Let $K$ be a convex parameter set of diameter $D$, and suppose the loss gradients over $K$ have norm at most $G$. OGD with step sizes $\eta_t=D/(G\sqrt t)$ satisfies
\begin{equation}
    \sum_{t=1}^{T} \ell_t(M_t)
    -
    \min_{M \in K}\sum_{t=1}^{T}\ell_t(M)
    \le
    \frac{3}{2}GD\sqrt T .
    \label{eq:ogd-regret}
\end{equation}
In particular, finite-memory linear predictors are learnable with $O(\sqrt T)$ regret, with constants depending on the memory length, norm radius, and data bounds.
\end{theorem}

There is also an exact structural reason to study linear predictors. In the noiseless case, every finite-dimensional LDS induces a finite-memory autoregressive predictor.

\begin{theorem}[Cayley--Hamilton representation; cf.\ {\citet[Thm.~9.4]{hazan2025onlinecontrol}}]
\label{thm:cayley-hamilton}
Consider the noiseless LDS
\[
    x_{t+1}=Ax_t+Bu_t,\qquad y_t=Cx_t,
\]
where $A\in\R^{d_x\times d_x}$ has characteristic polynomial
\[
    p_A(\lambda)=\lambda^{d_x}+a_1\lambda^{d_x-1}+\cdots+a_{d_x}.
\]
Then there exist matrices $M_{1,i}\in\R^{d_y\times d_y}$ and $M_{2,j}\in\R^{d_y\times d_u}$, for $i,j=1,\ldots,d_x$, such that for every $t>d_x$,
\begin{equation}
    y_t
    =
    \sum_{i=1}^{d_x} M_{1,i} y_{t-i}
    +
    \sum_{j=1}^{d_x} M_{2,j} u_{t-j}.
    \label{eq:cayley-hamilton-predictor}
\end{equation}
\end{theorem}

Theorem~\ref{thm:cayley-hamilton} shows that linear predictors are not merely a heuristic: in the noiseless setting they can exactly represent finite-dimensional LDS predictions after a transient. The drawback is that the memory length depends on the hidden dimension $d_x$. Lemma~\ref{lem:stable-truncation} removes this dependence for strictly stable systems, but replaces it with dependence on the stability margin. Spectral filtering addresses precisely this bottleneck.

%------------------------------------------------------------------------------
\subsection{Spectral Filtering: Compressing Long Memory}
\label{sec:spectral-filtering-prelim}
%------------------------------------------------------------------------------

Spectral filtering was introduced to learn LDSs without explicit system identification and without a long finite-memory window~\citep{hazan2017learning}. The exposition in \citet[Ch.~11]{hazan2025onlinecontrol} begins with the scalar impulse-response vectors
\begin{equation}
    \mu_\alpha
    =
    \begin{bmatrix}
    1 & \alpha & \alpha^2 & \cdots & \alpha^{T-1}
    \end{bmatrix}^{\!\top},
    \qquad \alpha \in [0,1],
    \label{eq:impulse-vector}
\end{equation}
and the Hankel/Hilbert matrix
\begin{equation}
    Z_T
    :=
    \int_0^1 \mu_\alpha\mu_\alpha^\top\,d\alpha
    \in \R^{T\times T},
    \qquad
    (Z_T)_{ij}=\frac{1}{i+j+1}
    \quad (i,j=0,\ldots,T-1).
    \label{eq:hankel-matrix}
\end{equation}
Let $\phi_1,\ldots,\phi_T$ be the eigenvectors of $Z_T$, ordered by decreasing eigenvalue. The eigenvalues of $Z_T$ decay rapidly; one convenient form is
\begin{equation}
    \sigma_k(Z_T) \le 4\pi^2 \exp\!\left(-\frac{k}{\log T}\right)
    \qquad (k,T>1),
    \label{eq:hankel-decay}
\end{equation}
as stated in \citet[Thm.~11.1]{hazan2025onlinecontrol}. This decay implies that the first polylogarithmically many eigenvectors approximate the family of impulse responses $\{\mu_\alpha:\alpha\in[0,1]\}$ well.

The resulting spectral predictor class uses the features $\langle \tilde u_t,\phi_i\rangle$, where
\[
    \tilde u_t=(u_{t-1},u_{t-2},\ldots,u_1,0,\ldots,0).
\]
In the scalar case,
\begin{equation}
    \hat y_t
    =
    \sum_{i=1}^{h} M_i\, \langle \tilde u_t,\phi_i\rangle.
    \label{eq:spectral-predictor}
\end{equation}
For symmetric LDSs with real spectrum, diagonalizing $A$ reduces the high-dimensional impulse response to a sum of scalar responses of the form~\eqref{eq:impulse-vector}. This yields an approximation whose feature count is polylogarithmic in $T$ and does not scale directly with the hidden dimension. The online learner then runs a convex online algorithm over the coefficients $M_i$.

\begin{remark}[What spectral filtering contributes]
Spectral filtering replaces a long raw-history basis by a compressed basis tailored to LDS impulse responses. It is the immediate mathematical predecessor of the present work. Our method does not learn Hankel eigenfeatures; instead, it injects one fixed polynomially filtered channel into a transformer. The shared idea is that a carefully chosen linear transformation of the input sequence can expose long-range dynamical structure to a downstream learner.
\end{remark}

%------------------------------------------------------------------------------
\subsection{Difficult Modes: Marginal Stability and Asymmetry}
\label{sec:difficult-modes}
%------------------------------------------------------------------------------

The preceding stable-system picture becomes less favorable when eigenvalues approach the unit circle. If $\rho(A)=1$, geometric truncation no longer gives a finite-memory approximation with rapidly vanishing tail. Moreover, for general asymmetric systems, the hidden dimension can be a genuine statistical barrier. For example, \citet[Thm.~12.1]{hazan2025onlinecontrol} gives a cyclic-permutation LDS construction in which any online learner incurs expected cumulative squared loss $\Omega(d_x)$ over $T\ge d_x$ rounds before the hidden state is revealed.

The same chapter describes Luenberger observers as one way to reshape difficult partially observed systems. If $(A,C)$ is observable, an observer gain $L$ can be chosen so that the error dynamics evolve under
\[
    e_{t+1} = (A-LC)e_t,
\]
and pole placement can move the spectrum of $A-LC$ into a strictly stable region. This converts a difficult partially observed prediction problem into a stable one, up to a burn-in term governed by the stability margin and eigenbasis conditioning.

Universal sequence preconditioning can be read in the same broad tradition: reshape the sequence or dynamics so that a downstream learner faces an easier spectral object.

%------------------------------------------------------------------------------
\subsection{Polynomial Preconditioning of LDS Observations}
\label{sec:polynomial-preconditioning-prelim}
%------------------------------------------------------------------------------

Let $B$ denote the backshift operator, $(Bx)_t=x_{t-1}$. For a polynomial
\[
    p(z)=\sum_{k=0}^{d}a_k z^k,
\]
define the filtered observation sequence
\begin{equation}
    \tilde y_t=(p(B)y)_t=\sum_{k=0}^{d} a_k y_{t-k}.
    \label{eq:poly-filter-output}
\end{equation}
The central algebraic identity behind universal sequence preconditioning is the following.

\begin{lemma}[Polynomial filtering as spectral reshaping; cf.\ {\citet{marsden2024universal}}]
\label{lem:polynomial-identity}
Consider the noise-free autonomous LDS $x_{t+1}=Ax_t$, $y_t=Cx_t=C A^t x_0$. Let
\[
    p^*(z):=z^d p(1/z)=\sum_{k=0}^{d} a_k z^{d-k}
\]
be the reciprocal polynomial. Then
\begin{equation}
    (p(B)y)_t
    =
    C\,p^*(A)\,A^{t-d}x_0.
    \label{eq:poly-identity}
\end{equation}
\end{lemma}

\begin{proof}
By direct calculation,
\[
    (p(B)y)_t
    =
    \sum_{k=0}^{d} a_k C A^{t-k}x_0
    =
    C\left(\sum_{k=0}^{d}a_k A^{d-k}\right)A^{t-d}x_0
    =
    C\,p^*(A)\,A^{t-d}x_0.
\]
\end{proof}

Thus filtering the observed sequence once is equivalent, algebraically, to inserting $p^*(A)$ into the hidden dynamics. If $A$ is diagonalizable with real spectrum in $[-\rho,\rho]$, then
\[
    \rho(p^*(A))
    \le
    \sup_{\lambda\in[-\rho,\rho]}\abs{p^*(\lambda)}.
\]
The right-hand side is a worst-case spectral upper bound over the interval. It is not a claim that a single polynomial is optimal for every fixed finite spectrum; rather, it is the uniform criterion used by Chebyshev approximation.

To minimize this worst-case bound, choose the spectral polynomial
\[
    q_d(\lambda)=\rho^d\,\tilde T_d(\lambda/\rho),
\]
where $\tilde T_d$ is the monic Chebyshev polynomial. The FIR polynomial applied to the sequence is then the reciprocal
\[
    p_d(z)=q_d^*(z):=z^d q_d(1/z),
\]
so that $p_d^*(A)=q_d(A)$ in Lemma~\ref{lem:polynomial-identity}.

\begin{theorem}[Chebyshev minimax]
\label{thm:chebyshev-minimax-prelim}
Among all real monic polynomials $q$ of degree $d\ge1$, the monic Chebyshev polynomial $\tilde T_d$ uniquely minimizes the uniform norm on $[-1,1]$:
\[
    \max_{x\in[-1,1]}\abs{\tilde T_d(x)}
    =
    \min_{\substack{q\ \mathrm{monic}\\ \deg(q)=d}}
    \max_{x\in[-1,1]}\abs{q(x)}
    =
    2^{1-d}.
\]
Consequently,
\[
    \sup_{\lambda\in[-\rho,\rho]}\abs{\rho^d\tilde T_d(\lambda/\rho)}
    =
    \rho^d 2^{1-d}.
\]
\end{theorem}

This theorem is the precise mathematical reason Chebyshev coefficients are a natural default: they minimize the worst-case spectral amplification among all monic degree-$d$ polynomials on a real spectral interval. In the regret bounds of \citet{marsden2024universal}, this spectral term trades off against a noise-gain term depending on the coefficient norm of $p_d$. We use this result as an inductive-bias principle, not as a theorem about transformers.

%------------------------------------------------------------------------------
\subsection{From Polynomial Preconditioning to Patch Transformers}
\label{sec:patch-transformer-connection}
%------------------------------------------------------------------------------

Patch-based time-series transformers, including PatchTST and Moirai-style models~\citep{nie2023patchtst,liu2024moirai2}, split a sequence into non-overlapping patches of length $P$ and embed each patch independently:
\[
    \mathbf p_i=(x_{(i-1)P+1},\ldots,x_{iP}),
    \qquad
    \mathbf e_i=W\mathbf p_i+b.
\]
Before the first attention layer, $\mathbf e_i$ contains information only from its own patch. Temporal structure spanning several patches must be reconstructed by attention from scratch. This is the cross-patch information bottleneck.

Our method injects a polynomially filtered residual channel before patching. Let $s=P$ be the stride and let $p_d$ be the reciprocal Chebyshev FIR polynomial above. The full filtered signal is
\[
    h_t=(p_d(B_s)x)_t,
    \qquad
    B_s x_t=x_{t-s}.
\]
Because the raw target channel $x_t$ is already provided to the transformer, the auxiliary channel is the residual
\begin{equation}
    r_t=h_t-x_t=(p_d(B_s)-1)x_t.
    \label{eq:residual-channel}
\end{equation}
For the main configuration $d=4$ and $s=P=16$,
\[
    p_4(z)=1-z^2+\frac18 z^4,
    \qquad
    r_t=-x_{t-32}+\frac18 x_{t-64}.
\]
Thus each current patch receives, through its widened input projection, fixed information about observations two and four patches in the past. The transformer still predicts the original target, not the preconditioned target.

\begin{remark}[Scope of the theory]
The LDS and online-learning results above do not prove that a transformer trained on heterogeneous real-world time series must improve. They justify a design principle: low-degree Chebyshev filters are minimax spectral shapers for real LDS modes; spectral and polynomial preprocessing can reduce the burden on a downstream learner; and patch transformers have an architectural bottleneck exactly where such preprocessing can help. The empirical sections of the paper test whether this inductive bias is useful in practice.
\end{remark}

\begin{thebibliography}{99}

\bibitem[Davis(1975)]{davis1975interpolation}
P.~J. Davis.
\newblock \emph{Interpolation and Approximation}.
\newblock Dover Publications, 1975.

\bibitem[Hazan and Singh(2025)]{hazan2025onlinecontrol}
E.~Hazan and K.~Singh.
\newblock \emph{Introduction to Online Control}.
\newblock Version 1.1 beta, 2025.

\bibitem[Hazan et~al.(2017)Hazan, Singh, and Zhang]{hazan2017learning}
E.~Hazan, K.~Singh, and C.~Zhang.
\newblock Learning linear dynamical systems via spectral filtering.
\newblock In \emph{Advances in Neural Information Processing Systems}, 2017.

\bibitem[Liu et~al.(2024)]{liu2024moirai2}
Y.~Liu et~al.
\newblock Moirai-2: A universal time series foundation model.
\newblock 2024.

\bibitem[Marsden and Hazan(2025)]{marsden2024universal}
A.~Marsden and E.~Hazan.
\newblock Universal sequence preconditioning.
\newblock arXiv:2502.06545, 2025.

\bibitem[Nie et~al.(2023)]{nie2023patchtst}
Y.~Nie, N.~H. Nguyen, P.~Sinthong, and J.~Kalagnanam.
\newblock A time series is worth 64 words: Long-term forecasting with transformers.
\newblock In \emph{International Conference on Learning Representations}, 2023.

\end{thebibliography}

\end{document}
